\subsection{有限性定理}

引理 7. 设 A 是 T 中的单位向量集。若存在实数 \(\lambda < 1\)，使得 \((a|a') \leqslant \lambda\) 对于 \(a, a' \in \mathbf{A}\) 且 \(a \neq a'\) 成立，则 A 是有限集。

对于 \(a, a' \in A\) 且 \(a \neq a'\)，有

\[
\left\| a - a ^ {\prime} \right\| ^ {2} = 2 - 2 (a \mid a ^ {\prime}) \geqslant 2 - 2 \lambda .
\]

现在，T 的单位球面 S 是紧的，因此存在一个由直径 \(<(2-2\lambda)^{1/2}\) 的集合组成的有限覆盖，并且这些集合中每个至多包含 A 的一个点，于是引理得证。

记 \(U(w)\) 为与从 E 到自身的仿射映射 \(w \in W\) 相关联的 T 的自同构。我们有

\[
w (x + t) = w (x) + \operatorname{U} (w). t \text {for} t \in \mathrm{T} \text {and} x \in \mathrm{E}.
\]

这定义了从群 W 到 T 的正交群的同态 U；U 的核是属于 W 的平移。

定理 3. (i) 一个室的墙的集合是有限的。

\begin{enumerate}
\def\labelenumi{(\roman{enumi})}
\setcounter{enumi}{1}
\item
属于 \(\mathfrak{H}\) 的超平面的方向的集合是有限的。
\item
群 U(W) 是有限的。
\end{enumerate}

断言 (i) 立即由命题 3(iii) 和引理 7 得出。

设 \(\mathbf{C}\) 是一个室，\(\mathfrak{M}\) 是其墙的集合。\(\overline{\mathbf{C}}\) 的面片（相对于 \(\mathfrak{H}\)）与相对于 \(\mathfrak{M}\) 的相同（§ 1，第 4 号，命题 9）。

由于 M 是有限的，面片的数目也是有限的。由于一个面片只与有限个属于 H 的超平面相交，所以与 \(\overline{C}\) 相交的属于 H 的超平面集合是有限的，因而 A(C)（即与某个属于 H 且与 \(\overline{C}\) 相交的超平面正交的 T 中单位向量的集合）也是有限的。因此，存在实数 \(\lambda < 1\)，使得 \((a|a') \leqslant \lambda\) 对于 \(a, a' \in A(C)\) 且 \(a \neq a'\) 成立。

令 A 为与属于 H 的某个超平面正交的 T 中单位向量所成的集合。设 \(a, a' \in A\) 且 \(a \neq a'\)。若 a 与 \(a'\) 平行，则 \(a = -a'\) 且 \((a|a') = -1\)。否则，令 \(H \in H\)（分别为 \(H' \in H\)）满足 a（分别为 \(a'\)）与 H（分别与 \(H'\)）正交。我们有 \(H \cap H' \neq \varnothing\)，且若 \(x \in H \cap H'\)，则存在元素 \(w \in W\)，使得 \(x \in w(\overline{C})\)。于是，向量 \(U(w).a\) 和 \(U(w).a'\) 属于 A(C)，且有

\[
(a \mid a ^ {\prime}) = (\mathrm{U} (w). a \mid \mathrm{U} (w). a ^ {\prime}) \leqslant \lambda ,
\]

且由引理 7，集合 \(\mathbf{A}\) 是有限的。因此 (ii)。

现在，设 \(w \in W\)，满足 \(\mathrm{U}(w).a = a\) 对所有 \(a \in A\) 成立。则 \(\mathrm{U}(w).t = t\) 对 T 中由 A 生成的子空间内的所有 t 成立。另一方面，若 \(t \in T\) 与 A 正交，则对所有 \(\mathrm{U}(s_{\mathrm{H}}).t = t\) 都有 \(H \in H\)，因而 \(\mathrm{U}(w).t = t\)，最终 \(\mathrm{U}(w) = 1\)。由于对所有 \(\mathrm{U}(w)(\mathrm{A}) = \mathrm{A}\) 都有 \(w \in W\)，我们得到 \(\mathrm{U}(\mathrm{W})\) 同构于有限集 A 的一个置换群，因此 (iii)。

命题 4. 设 C 是一个室，N 是 C 的一组墙。令 \(W_{N}\) 为由关于 N 中元素的正交反射生成的 W 的子群。对于 \(H \in N\)，记 \(e_{H}\) 为与 H 正交且位于 H 与 C 同一侧的单位向量。下列条件等价：

\begin{enumerate}
\def\labelenumi{(\roman{enumi})}
\item
群 \(\mathbf{W}_{\mathfrak{N}}\) 是有限的。
\item
存在 \(\mathbf{E}\) 中一点，在 \(W_{\mathfrak{N}}\) 的每个元素下不变。
\item
属于 \(\mathfrak{N}\) 的超平面有非空交。
\item
向量族 \((e_{\mathrm{H}})_{\mathrm{H}\in \mathfrak{N}}\) 在 T 中自由。
\end{enumerate}

根据 §3 开头的性质 (D2)，W 中每一点的稳定子都是有限的，因此 (ii) 蕴含 (i)。

由于群 \(W_{N}\) 由关于属于 N 的超平面所成的反射集合生成，\(W_{N}\) 的不动点就是属于每个超平面 \(H \in N\) 的 E 中的点，因此 (ii) 与 (iii) 等价。

假设 \(a\) 是 \(E\) 中的一点，使得 \(a \in H\) 对所有 \(H \in \mathfrak{N}\) 成立，并令 \(t \in T\) 满足 \(a + t \in C\)。由于 \((e_H | e_{H'}) \leqslant 0\) 对满足 \(H, H' \in \mathfrak{N}\) 且 \(H \neq H'\) 的情形成立（命题 3），又由于 \((t | e_H) > 0\) 对所有 \(H \in \mathfrak{N}\) 成立，引理 3(ii) 蕴含 \(e_H\)（其中 \(H \in \mathfrak{N}\)）线性无关。因此，(iii) 蕴含 (iv)。

最后，假设向量族 \((e_{\mathrm{H}})_{\mathrm{H}\in\mathfrak{N}}\) 是自由的。令 a 为 E 中一点。对于任意 \(H\in N\)，存在实数 \(c_{H}\)，使得 H 恰好由 E 中满足 \(a+t\) 的点 \((t|e_{\mathrm{H}})=c_{\mathrm{H}}\) 组成。由于族 \((e_{\mathrm{H}})\) 是自由的，存在 \(t\in T\)，使得 \((t|e_{\mathrm{H}})=c_{\mathrm{H}}\) 对所有 \(H\in N\) 成立，且 E 中的点 \(a+t\) 属于所有超平面 \(H\in N\)。因此，(iv) 蕴含 (iii)。

注. 1) 由于 W 由关于室 C 的墙的反射生成，前一命题给出了 W 有限的一个判据。我们将在第 9 号回到这个问题。

\begin{enumerate}
\def\labelenumi{\arabic{enumi})}
\setcounter{enumi}{1}
\tightlist
\item
设 F 是有限维实仿射空间，G 是 F 的自同构群。对于所有 \(g \in G\)，记 \(U(g)\) 为与 g 相关的 F 的平移空间 V 的自同构。假设 \(U(G)\) 的像是 \(\mathbf{GL}(V)\) 的有限子群；则 V 具有一个在 \(U(G)\) 下不变的标量积（《积分》第 VII 章 §3 第 1 号命题 1）。此外，如果 G 配备离散拓扑后在 F 上适当地作用，且由反射生成，则可以将本段的结果应用于 G。
\end{enumerate}

